\section{Controlled model approximation changes three parts of inference}
\label{sec:approximation-control}

Correctness of the target and valid sources alone is insufficient for the
previous leave-out dispersion proof: the variance correction also involves
invalid sources. We make the missing terms explicit. The controlled extension
below requires declared or separately certified error budgets. It does not
estimate arbitrary nonlinear approximation errors for free. If budgets are
estimated rather than declared, their certificate-failure probabilities must
be included separately in the total error allocation. The present corrected
experiments use fixed declared envelopes.

For one arm or a stacked contrast, retain $H_j$, $D_j=\diag(d_j)$ and the
quadratic matrix $Q$ from the linear branch. Here $n_j$ counts observations
in the current arm; contrast bounds combine the two arm-specific quantities. Let
\[
 m_j=\E(W_j\mid\mathcal D),\qquad u_j=(I-H_j)m_j,\qquad
 \zeta_j=m_j^TD_ju_j.
\]
Direct expectation gives
\[
 \E(\widehat v_j\mid\mathcal D)=v_j+\zeta_j,\qquad
 \E(\widehat D\mid\mathcal D)=D-a_K\sum_j\zeta_j.
\]
In particular, positive $\zeta_j$ makes the correction subtract too much noise and
can hide source dispersion. A budget only for valid-source mean bias does not
repair this error.

\subsection{Observable design factors and declared residual envelopes}

Suppose $\|u_j\|_2\le\rho_j$. For a uniform approximation error
$\|r_j\|_\infty\le\epsilon_j$ in
$\E(Y_j\mid\mathcal D)=X_j\beta_j+r_j$, one may use
$\rho_j=\epsilon_j\sqrt{n_j}$, because $I-H_j$ is a contraction.
Let $d_j^{\max}=\max_i d_{ji}$, $G_j=\sum_i d_{ji}^2$ and
$c_j=\|(I-H_j)D_jH_j\|_{\rm op}$. Then
\[
 |\zeta_j|\le B_j:=\min\left\{
   \tfrac12\sqrt{G_j}\rho_j,
   d_j^{\max}\rho_j^2+\tfrac12\sqrt{n_j}\,c_j\rho_j\right\}.
\]
The first inequality uses $\|m_j\|_\infty\le1/2$. For the second, write
$m_j=H_jm_j+u_j$ and use $\|H_jm_j\|_2\le\sqrt{n_j}/2$.
With thin QR basis $U_j$,
\[
 c_j^2=\lambda_{\max}\left\{
 U_j^TD_j^2U_j-(U_j^TD_jU_j)^2\right\}.
\]
Thus the design factors are computable without a dense patient matrix.

If at most $q$ sources can be nonlinear, define
\[
 C_B=\operatorname{TopSum}_q(B_j),\qquad
 C_U=\operatorname{TopSum}_q\{(d_j^{\max})^2\rho_j^2\},
\]
where $\operatorname{TopSum}_q$ sums the $q$ largest entries. This protects an
unknown nonlinear subset. If no count restriction is justified, use $q=K$.
For a contrast, construct these sums separately by arm and add them; no common
nonlinear subset is assumed. In the targeted experiment, target and valid
sources remain linear, so $q_a\le K-g_a$ is justified. The treated-arm uniform
approximation envelope is fixed in advance across all simulated amplitudes.

\subsection{Correct the quadratic certificate, not just its center}

The patient matrix remains zero diagonal, but its linear term changes to
\[
 (Qm)_j=\gamma_j(A\Gamma^Tm)_j-\tfrac{a_K}{2}
        \{D_ju_j+(I-H_j)D_jm_j\}.
\]
Therefore
\[
 \|Qm\|_2^2\le LD+
  \frac{a_K^2}{2}\left(\sqrt{C_U}+\tfrac12\sqrt{\sum_jG_j}\right)^2.
\]
Let $Q_0$ be the exact-linear constant. Replace it by
\[
 Q_0^{\rm app}=Q_0+\frac{a_K^2}{2}
          \left(C_U+\sqrt{C_U\sum_jG_j}\right).
\]
The same Gaussian MGF domination now applies to
$\widehat D-\E(\widehat D\mid\mathcal D)$, with variance proxy
$LD+Q_0^{\rm app}$ and the same scale $c$. Since
$a_K\sum_j\zeta_j\le a_K C_B$, use
\[
 y=\widehat D+a_K C_B+c\log(1/\alpha_D)
\]
in the earlier quadratic inversion, together with $Q_0^{\rm app}$.
This gives a valid upper bound for $D$. At zero projection error it reduces
exactly to the preceding linear-model certificate.

\subsection{Arbitrary bounded invalid means require no supplied amplitude}

If the target and the guaranteed valid sources are exactly linear, only the
remaining at most $K-g_a$ sources per arm can have nonzero projection residuals.
For those sources, bounded outcomes alone imply
\[
 \|u_{ja}\|_2\le\|m_{ja}\|_2\le\tfrac12\sqrt{n_{ja}}.
\]
Thus the preceding formulas can use $\rho_{ja}=\sqrt{n_{ja}}/2$ and
$q_a=K-g_a$ without knowing an invalid source's nonlinearity amplitude.
The valid-source mean-bias allowance remains zero. This gives conditional
coverage with arbitrary bounded invalid-source means, under the retained
linear target/valid-source and valid-count assumptions.

The smaller declared-envelope option uses a narrower model class and may be
more precise. The universal bounded-invalid option pays for a larger possible
variance-centering error. Its present bound can have an $O(n_s^{-1/2})$
source-radius floor as $K$ grows, even in favorable data. This is a limitation
of this certificate, not a proved minimax lower bound. The exact-linear
fourth-root rate must not be attributed automatically to the broader option.

\subsection{Mean-bias budgets remain necessary when valid models are approximate}

If a valid source's conditional score mean differs from $\mu_{t,X}^a$ by at
most $h_{ja}$, its unknown valid-subset mean is bounded by the average of the
$g_a$ largest $h_{ja}$. Add those arm-specific averages to the arm dispersion
radius. For the contrast certificate, average the $g_\cap$ largest
$h_{j1}+h_{j0}$. The target conditional band likewise needs its own mean-bias
allowance.

For a common transported mean of the form $\phi^T\beta_t^a+r_t^a$, with
$\sup_x|r_t^a(x)|\le\epsilon_t^a$, exact balance supplies the computable
valid-source allowance
\[
 h_{ja}=\epsilon_t^a(1+\|\gamma_{ja}\|_1),
\]
and the target allowance is the analogous sum over its two arms. These bounds
are conservative. For an approximately linear target, use universal composition
protection unless the target coefficient certificate is also extended to cover
its approximation bias. The empirical Bernstein target experiment in this
version retains an exactly linear target.

\subsection{A conditional counterexample and what it establishes}

Fix each treatment arm's design to $N/4$ observations at $x=-1$, $N/2$ at
$x=0$, and $N/4$ at $x=1$, with $N=500$. Fit the basis $(1,x)$ and balance
to mean zero; then $\gamma_i=1/N$. Valid treated means are $.5$, control means
are $.4$, and one quarter of sources have treated means
\[
 .5+\delta_K+\eta f(x),\qquad f(x)=x^2-.5,\quad
 \delta_K=2\sqrt{.5/1000}\,K^{-1/4}.
\]
All target and valid-source models are linear. For fixed $\eta>0$, $f$ is
orthogonal to the fitted basis, while the invalid-source average shift is
$\delta_K$. Thus true source dispersion is
$D=\tfrac14\tfrac34\delta_K^2\to0$. But
\[
 \zeta_{\rm invalid}
 =\sum_g N_g d_g\{\delta_K+\eta f(x_g)\}\eta f(x_g)
 \longrightarrow \eta^2\sum_gN_gd_g f(x_g)^2>0.
\]
Consequently, the uncorrected dispersion statistic converges to a negative
constant as $K$ increases with this fixed design and $N$. Its nominal upper
bound eventually becomes zero, while the pooled bias remains larger than its
$K^{-1/2}$ outcome-noise scale. The source conditional interval can therefore
lose coverage. A wide population enlargement can conceal that failure.

This is a counterexample to the proposed \emph{conditional} extension without
model-error correction. It is not presented as an i.i.d. random-design
population simulation. The exact grouped-Bernoulli experiment and the separate
random-design study are reported distinctly. Neither good population coverage
nor this counterexample alone proves a general nonlinear procedure.
